\begin{answers}{ch:tokenizer}{トークン化：文字列を数の列にする}
\item 列は 8 個から 4 個になり、「いない」「いない」「ば」「あ」です。
最初は（い, な）と（な, い）が 2 回で並ぶので、先に現れた（い, な）を選び、「いな」を作ります。
列は「いな」「い」「いな」「い」「ば」「あ」の 6 個になります。
数え直すと（いな, い）が 2 回で最多なので、これをまとめて「いない」を作ります。
\item 答えは \code{e7 8c ab}（10 進法で 231, 140, 171）です。
16 進法の 7, 3, 2, B を 2 進数にすると、\code{0111 0011 0010 1011} の 16 桁です。
これを 4・6・6 桁に分けると \code{0111}、\code{001100}、\code{101011} です。
目印を付けて \code{1110}\,\code{0111}、\code{10}\,\code{001100}、\code{10}\,\code{101011} とし、8 桁ずつ 16 進法で読むと \code{e7}、\code{8c}、\code{ab} です。
\code{list("猫".encode("utf-8"))} も \code{[231, 140, 171]} を返し、一致します。
\item 9 個のトークン \code{<e3 81> <a7> <e3 81> <82> <e3> <82> <8b> <e3 80> <82>} になり、1 個にまとまりません（Python で確かめました）。
逆の順では、規則 10・9・8・7・6 は、材料の「である」「。」「\code{<e3 80>}」などがまだ無いので、何も起きません。
規則 5 で「\code{e3 80}」が一つできますが、規則 4・3・2 も材料が無く空振りします。
最後の規則 1 で「\code{e3 81}」が二つできて終わります。
後から学習した規則は、先に学習した規則でできたトークンを材料にすることがあるので、学習した順に当てはめないと、長いトークンができません。
\item 先頭から 3 桁ずつ区切られ、余りの 1〜3 桁が最後に残ります（実測）。
「\code{123456789012}」の先頭 $k$ 桁（$k=1,\dots,12$）を切ると、$k\le3$ は 1 トークンです。
4 桁は \code{123}/\code{4}、8 桁は \code{123}/\code{456}/\code{78}、10 桁は \code{123}/\code{456}/\code{789}/\code{0} です。
トークン数は、$k$ を 3 で割って切り上げた数です。
区切りは\textbf{左から}なので、位取りの 3 桁区切り（右から）とは合いません。
たとえば「\code{1000000}」は \code{100}/\code{000}/\code{0} になります。
\item 列は、足したトークンの数のわりには縮まず、トークン 1 種類あたりの学習例は激減し、パラメータは約 8 倍に増えるはずです（10 万は実測していません）。
列の長さは、\tabref{tab:tok-tradeoff}のとおり、語彙を 2 倍にするごとの伸びは 512 から先ほぼ一定（0.25〜0.3 文字）ですが、足したトークン 100 個あたりの伸びは 0.21、0.05、0.025、0.014 文字と小さくなっています。
最後のマージ（3,832 番目）でも組は 85 回しか現れないので（\figref{fig:tok-tradeoff}の右）、ここから 9 万個あまり足しても、1 個あたりの伸びはさらに小さくなります。
学習例は、\figref{fig:tok-vocab-cost}の左のとおり、語彙 4,088 でもすでに 1 種類あたり中央値 164 回まで減っています。
増えるトークンの多くは、この小説に数回しか現れない並びです。学習例が少ないので、その埋め込みはうまく学べません。
パラメータは、埋め込み表だけで $100{,}000\times192=1{,}920$ 万個です。
埋め込み以外の部分（約 260 万 $-$ 79 万 $\approx$ 181 万個）と合わせて約 2,100 万個で、9 割が埋め込みです。
語彙を増やすなら、学習に使う文章も増やす必要があります。
\end{answers}
